% Work/MAT -- hodina 29; druha a treti odmocnina BEZ racionalnich exponentu.
% Podklady: ucebnice Mocniny a Odmocniny.tex, oddil Odmocniny;
% Realisticky.cz 010207 Druha odmocnina, 010208 Treti odmocnina.
% Presne hodnoty, geometricka motivace, podminky, omyly, vsechna reseni vlastni.
\documentclass[11pt,a4paper]{article}
% Spolecny styl pro ucitelske pripravy MAT (Work/MAT)
% Verze 1.1, 2026-10-23
\NeedsTeXFormat{LaTeX2e}
\ProvidesPackage{priprava}[2026/10/23 v1.1 Sablona ucitelskych priprav MAT]

% Zakladni typografie a cestina
\RequirePackage[T1]{fontenc}
\RequirePackage[utf8]{inputenc}
\RequirePackage[czech]{babel}
\RequirePackage{lmodern}
\RequirePackage{microtype}
\RequirePackage{amsmath,amssymb}
\RequirePackage{xcolor}
\RequirePackage{enumitem}
\RequirePackage[most]{tcolorbox}
\RequirePackage{tabularx,booktabs,array}
\RequirePackage{fancyhdr}
\RequirePackage{needspace}
\RequirePackage{geometry}
\RequirePackage[hidelinks]{hyperref}

\geometry{a4paper,top=17mm,bottom=18mm,left=23mm,right=23mm}
\definecolor{notegray}{gray}{0.39}
\setlength{\parindent}{0pt}
\setlength{\parskip}{2.5pt plus 1pt minus 0.5pt}
\setlist[itemize]{topsep=2pt,itemsep=1pt,parsep=0pt,leftmargin=1.45em}
\setlist[enumerate]{topsep=2pt,itemsep=2pt,parsep=0pt,leftmargin=1.65em}
\setcounter{secnumdepth}{0}
\renewcommand{\arraystretch}{1.12}

% Zahlavi: vlevo rocnik | globalni tema, vpravo cislo: nazev hodiny
% Zapati zustava zcela prazdne.
\pagestyle{fancy}
\fancyhf{}
\renewcommand{\headrulewidth}{0.35pt}
\renewcommand{\footrulewidth}{0pt}
\setlength{\headheight}{14pt}
\newcommand{\nastavPripravu}[4]{%
  \fancyhead[L]{\small #1\enspace\textbar\enspace #2}%
  \fancyhead[R]{\small #3: #4}%
}

% Minimalisticke cernobile boxy. Kazda definice zvlast.
\tcbset{commonlesson/.style={enhanced,breakable=false,colback=white,colframe=black,
  colbacktitle=white,coltitle=black,boxrule=0.45pt,
  arc=0.8mm,outer arc=0.8mm,left=3mm,right=3mm,top=1.6mm,bottom=1.6mm,
  toptitle=1mm,bottomtitle=0.7mm,fonttitle=\bfseries,
  attach boxed title to top left={xshift=2.5mm,yshift=-1.15mm},
  boxed title style={colback=white,colframe=white,boxrule=0pt,left=1mm,right=1mm}}}
\newtcolorbox{zadaniBox}[1][]{commonlesson,title={Zadání},#1}
\newtcolorbox{prikladBox}[1][]{commonlesson,title={Příklad},#1}
\newtcolorbox{definiceBox}[1][]{commonlesson,title={Definice},#1}

% Spolecny box pro zadani a jeho reseni, bez mezery mezi castmi.
% Pouziti:
%   \begin{prikladReseniBox}
%     Zadani ulohy.
%     \tcblower
%     \begin{align*} ... \end{align*}
%   \end{prikladReseniBox}
% Spodni cast ma jemny sedy podklad a delici linku; lze v ni sazet
% vice radku matematiky. Obe casti maji jen jeden spolecny ramecek.
\tcbset{solutionlesson/.style={
  commonlesson,
  enhanced jigsaw,
  breakable,
  colbacklower=black!3,
  segmentation style={solid,draw=black!25,line width=0.35pt},
  fontlower=\small\color{notegray},
  before lower={\noindent\textbf{Řešení:}\enspace},
  leftlower=3mm,rightlower=3mm
}}
\newtcolorbox{prikladReseniBox}[1][]{solutionlesson,title={Příklad},#1}
\newtcolorbox{zadaniReseniBox}[1][]{solutionlesson,title={Zadání},#1}

% Pomucky pro autora pripravy: sedive poznamky; \reseni pro starsi pripravy
\newcommand{\poznamka}[1]{\par\begingroup\small\color{notegray}\itshape
  \textnormal{\textbf{Poznámka pro učitele.}}\enspace #1\par\endgroup}
\newcommand{\reseni}[1]{\par\begingroup\small\color{notegray}
  \textbf{Očekávané řešení:}\enspace #1\par\endgroup}
\newcommand{\blok}[1]{\needspace{5\baselineskip}\subsection*{#1}}
\newcommand{\vyrok}[1]{\uv{#1}}

\usepackage{tikz}
\nastavPripravu{1. ročník}{Mocniny a odmocniny}{29}{Druhá a třetí odmocnina}
\begin{document}
\poznamka{\textbf{45 minut:} 5 min motivace obsah/objem;
9 min definice druhé odmocniny + Příklad 1;
7 min Samostatná práce 1;
8 min definice třetí odmocniny + Příklad 2;
8 min Samostatná práce 2; 5 min nejčastější omyly a Příklad 3;
3 min shrnutí. Zbývající úlohy jsou rezerva.
\textbf{Dnes nepoužívej zápis mocnin se zlomkovým exponentem}; spojení až v hodině 30.}

\blok{Opačná otázka k umocňování: jaká byla délka strany?}
Čtverec se stranou $5\,\mathrm{cm}$ má obsah $25\,\mathrm{cm}^2$.
Když naopak známe \emph{obsah} $25\,\mathrm{cm}^2$, jak zjistíme délku strany?
U krychle podobně: objem $27\,\mathrm{cm}^3$ znamená hranu délky $3\,\mathrm{cm}$.
\begin{center}
\begin{tikzpicture}[scale=.72,line join=round]
\draw[thick] (0,0) rectangle (2.1,2.1);
\node at (1.05,1.1) {$S=25$};
\node[below] at (1.05,0) {$a=?$};
\begin{scope}[shift={(6.0,.25)}]
  \draw[thick] (0,0)--(2.0,0)--(2.0,1.75)--(0,1.75)--cycle;
  \draw[thick] (0,1.75)--(.75,2.4)--(2.75,2.4)--(2.0,1.75);
  \draw[thick] (2,0)--(2.75,.65)--(2.75,2.4);
  \node at (1,1) {$V=27$};
  \node[below] at (1,0) {$a=?$};
\end{scope}
\end{tikzpicture}
\end{center}
\poznamka{Jednotky u obrázků: první obsah je v $\mathrm{cm}^2$, druhý objem v $\mathrm{cm}^3$.
Žáci mohou odhadnout výsledné délky i před zavedením značky odmocniny.}

\blok{Druhá odmocnina: vždy volíme nezáporný výsledek}
\begin{definiceBox}[title={Definice: druhá odmocnina}]
Pro reálné $a\ge0$ je \textbf{druhá odmocnina} $\sqrt a$ jediné
\textbf{nezáporné} číslo $b$, pro které $b^2=a$.
\[
\boxed{\sqrt a=b\quad\Longleftrightarrow\quad b\ge0\ \text{a}\ b^2=a.}
\]
V oboru reálných čísel nemá $\sqrt a$ smysl pro $a<0$.
\end{definiceBox}
\begin{prikladReseniBox}[title={Příklad 1: co je pod odmocnítkem?}]
Urči $\sqrt{81}$, $\sqrt{0}$, $\sqrt{\frac{49}{100}}$,
$\sqrt{(-6)^2}$ a $-\sqrt{36}$.
\tcblower
$9$, $0$, $\frac7{10}$, $6$, $-6$.
Ve čtvrtém případě je $(-6)^2=36$, ale \textbf{druhá odmocnina z 36 je $+6$}.
V posledním příkladu je minus \emph{až před} odmocninou.
\end{prikladReseniBox}
\begin{zadaniReseniBox}[title={Samostatná práce 1: druhé odmocniny}]
\begin{enumerate}[label=\alph*)]
\item $\sqrt{144}$; \item $\sqrt{0{,}04}$;
\item $\sqrt{\frac9{25}}$; \item $\sqrt{(-8)^2}$;
\item $\sqrt{-16}$ v $\mathbb R$; \item $\sqrt{2^6}$;
\item $\sqrt{49+32}$; \item $\sqrt{36}+\sqrt{64}$.
\end{enumerate}
\tcblower
 a) $12$; b) $0{,}2$; c) $3/5$; d) $8$;
 e) není definována v $\mathbb R$; f) $8$;
 g) $\sqrt{81}=9$; h) $6+8=14$.
\end{zadaniReseniBox}

\blok{Třetí odmocnina: záporná čísla jsou povolena}
\begin{definiceBox}[title={Definice: třetí odmocnina}]
Pro libovolné $a\in\mathbb R$ je \textbf{třetí odmocnina} $\sqrt[3]{a}$
jediné reálné číslo $b$, pro které $b^3=a$:
\[
\boxed{\sqrt[3]{a}=b\quad\Longleftrightarrow\quad b^3=a.}
\]
Výsledek může být záporný: $\sqrt[3]{-8}=-2$, protože $(-2)^3=-8$.
\end{definiceBox}
\begin{prikladReseniBox}[title={Příklad 2: i záporný základ má svou třetí odmocninu}]
Vypočti $\sqrt[3]{64}$, $\sqrt[3]{-125}$,
$\sqrt[3]{-\frac8{27}}$ a $\sqrt[3]{(-4)^3}$.
\tcblower
$4$, $-5$, $-\frac23$, $-4$.
Znaménko se u liché mocniny zachová, proto třetí odmocnina existuje
i pro záporná reálná čísla.
\end{prikladReseniBox}
\begin{zadaniReseniBox}[title={Samostatná práce 2: třetí odmocniny a rozlišování}]
\begin{enumerate}[label=\alph*)]
\item $\sqrt[3]{216}$; \item $\sqrt[3]{-216}$;
\item $\sqrt[3]{0{,}001}$; \item $\sqrt[3]{-\frac1{64}}$;
\item $\sqrt[3]{(-7)^3}$; \item $\sqrt{(-7)^2}$;
\item $\sqrt[3]{-8}+\sqrt{16}$;
\item $\sqrt[3]{8+19}$.
\end{enumerate}
\tcblower
 a) $6$; b) $-6$; c) $0{,}1$; d) $-1/4$;
 e) $-7$; f) $7$; g) $-2+4=2$; h) $3$.
\end{zadaniReseniBox}

\blok{Pozor: tři typické chyby}
\begin{prikladReseniBox}[title={Příklad 3: odmocnina není totéž co řešení rovnice}]
Rozhodni, kolik výsledků má každý zápis:
$\sqrt{25}$; $x^2=25$; $\sqrt[3]{-27}$.
\tcblower
$\sqrt{25}=5$ má \textbf{jedinou} hodnotu.
Rovnice $x^2=25$ má \textbf{dvě} řešení $x=5$ a $x=-5$.
$\sqrt[3]{-27}=-3$ má jedinou reálnou hodnotu.
U odmocniny vždy pracujeme s její \emph{definovanou} hodnotou, ne s dvojicí řešení rovnice.
\end{prikladReseniBox}
\begin{zadaniReseniBox}[title={Samostatná práce 3: oprav nesprávný postup}]
\begin{enumerate}[label=\alph*)]
\item $\sqrt{16}=\pm4$.
\item $\sqrt{9+16}=\sqrt9+\sqrt{16}$.
\item $\sqrt{a^2}=a$ pro všechna $a\in\mathbb R$.
\item $\sqrt[3]{-8}$ neexistuje, protože je číslo pod odmocnítkem záporné.
\end{enumerate}
\tcblower
 a) Správně $\sqrt{16}=4$; znak $\pm$ se objeví při řešení $x^2=16$.
 b) Vlevo $\sqrt{25}=5$, vpravo $3+4=7$: odmocnina se přes součet \emph{nerozděluje}.
 c) Správně $\sqrt{a^2}=|a|$, třeba pro $a=-3$ vychází $3$, ne $-3$.
 d) $\sqrt[3]{-8}=-2$; třetí odmocnina ze záporného čísla existuje.
\end{zadaniReseniBox}
\poznamka{Větu $\sqrt{a^2}=|a|$ zdůvodni nezáporností druhé odmocniny.
Nevysvětluj ji chybným mechanickým \uv{rušením} dvou operací.
Pravidla pro součin a podíl odmocnin, částečné odmocňování a usměrňování
nech na navazující hodiny 33--36.}

\blok{Rozšiřující příklady: odhad a práce s proměnnými}
\begin{zadaniReseniBox}[title={Rozšiřující příklad 1: mezi kterými celými čísly?}]
\begin{enumerate}[label=\alph*)]
\item Odhadni bez kalkulačky $\sqrt{20}$ mezi dvě po sobě jdoucí celá čísla.
\item Odhadni $\sqrt[3]{30}$ a $\sqrt[3]{-30}$ mezi dvě celá čísla.
\item Pro která reálná $t$ má výraz $\sqrt{t-5}$ smysl?
\end{enumerate}
\tcblower
 a) $4<\sqrt{20}<5$, protože $16<20<25$.
 b) $3<\sqrt[3]{30}<4$ a $-4<\sqrt[3]{-30}<-3$.
 c) $t\ge5$.
\end{zadaniReseniBox}
\begin{zadaniReseniBox}[title={Rozšiřující příklad 2: opatrně s rovnicí}]
\begin{enumerate}[label=\alph*)]
\item Vyřeš $x^2=49$ a porovnej s $\sqrt{49}$.
\item Urči $\sqrt{(x-2)^2}$ pro $x=0$ a $x=5$.
\item Vyřeš $x^3=-64$.
\end{enumerate}
\tcblower
 a) $x=\pm7$, ale $\sqrt{49}=7$;
 b) $|x-2|$: pro $x=0$ je $2$, pro $x=5$ je $3$;
 c) $x=-4$.
\end{zadaniReseniBox}

\blok{Domácí procvičení}
\begin{zadaniReseniBox}[title={Domácí úkol 1: druhé a třetí odmocniny}]
\begin{enumerate}[label=\alph*)]
\item $\sqrt{225}$ a $\sqrt{\frac{121}{169}}$.
\item $\sqrt[3]{-343}$ a $\sqrt[3]{0{,}008}$.
\item $\sqrt{(-11)^2}$ a $\sqrt[3]{(-11)^3}$.
\item Je $\sqrt{-25}$ reálné číslo? A $\sqrt[3]{-125}$?
\item Vysvětli, proč $\sqrt{4+9}\ne\sqrt4+\sqrt9$.
\end{enumerate}
\tcblower
 a) $15$, $11/13$; b) $-7$, $0{,}2$;
 c) $11$, $-11$; d) první ne, druhý je $-5$;
 e) vlevo $\sqrt{13}$, vpravo $5$; protože $\sqrt{13}\ne5$.
\end{zadaniReseniBox}
\end{document}
